\begin{helper}
Let \(N,n\) be natural numbers.  Let
\(x\in\mathbb R^{\operatorname{Fin}(N)}\) be a nonincreasing vector with
nonnegative integer entries and total sum \(n\).  Then for every
\(r\le N\), the sum of the first \(r\) entries of \(x\) is at least
\(\min(r,n)\):
\[
\min(r,n)\le \sum_{i<r}x_i .
\]
\end{helper}
\begin{proof}
The case \(r=0\) is immediate, since both sides are \(0\).  Assume
\(0<r\le N\).  Because every \(x_i\) is a nonnegative integer, each
coordinate is either \(0\) or at least \(1\).

First suppose \(r\le n\).  If the first \(r\) entries had sum strictly less
than \(r\), then \(r\) nonnegative integer summands would have total less
than \(r\).  If the \(r\)-th entry were at least \(1\), then, because the
vector is nonincreasing, all first \(r\) entries would be at least \(1\), and
their sum would be at least \(r\), a contradiction.  Thus the \(r\)-th entry
is \(0\).  Nonnegativity and monotonicity then force every later entry to be
\(0\) as well.  Hence the total sum of all \(N\) entries is the same as the
sum of the first \(r\) entries, which is strictly less than \(r\le n\).  This
contradicts the total sum \(n\).  Therefore the first \(r\) entries have sum
at least \(r\), which is \(\min(r,n)\) in this case.

Now suppose \(n\le r\).  If some entry after the first \(r\) positions were
positive, then, because the vector is nonincreasing and has nonnegative
integer entries, all of the first \(r+1\) entries would be at least \(1\).
Their sum would then be at least \(r+1\), and hence greater than or equal to
\(n+1\), contradicting the total sum \(n\).  Thus every positive entry is
among the first \(r\) coordinates, so the first \(r\) entries contain the
entire total sum \(n\).  Hence their sum is \(n=\min(r,n)\).
\end{proof}
